\section{Discussion}
\label{sec:discussion}

\paragraph{Shrink towards a common model, not towards zero.}
The $\ell_{2,1}$ and trace-norm penalties share information between tasks, but their shrinkage target is the zero matrix (\cref{prop:limits}). When data are scarce, cross-validation chooses strong regularization, and the estimate of the part all schools have in common is shrunk along with everything else. Pooled ridge makes the opposite error: it estimates the common part with all the data and ignores school-specific effects. Partial pooling estimates the common part almost without penalty and spends the regularization budget only on deviations from it. On School this is enough to match the pooled model when schools have a handful of students and to beat every alternative once they have more. A trace-norm model can represent a shared mean as one of its directions, so the difference is not one of expressiveness but of where the penalty is applied. The behaviour of the partial-pooling estimators, close to pooling with little data and ahead of it with more, is what the adaptivity results of \citet{duan2023adaptive} predict for estimators that shrink towards a learned center.

\paragraph{Relation to earlier School results.}
At $\rho = 75\%$ our results agree with \citet{argyriou2008convex}: the pooled model explains 22.0\% of within-school variance against their 22.7\%, clearly below trace-norm MTL. They disagree with \citet{evgeniou2005kernel}, who found a single model for all schools best under the same protocol. Their models were support vector regressions rather than ridge regressions, and we did not reproduce their setup, so we cannot say which difference explains the disagreement. Neither study examined small training fractions, which is where, in our experiments, the pooled model is strongest and the choice between shrinking towards zero and shrinking towards a common model matters most.

\paragraph{Structure in the deviations.}
Constraining the deviations to be low rank does not help on School: from 10\% training onwards the two partial-pooling estimators are within 0.8\% of each other, and at 5\% the low-rank variant is worse (\cref{tab:main}). The rank that cross-validation selects for the deviations grows with the data (\cref{fig:structure}b), which suggests that schools differ along a few directions beyond their intercept, but these differences carry little predictive weight relative to the shared model and the noise. We expect structured deviations to matter more where task differences are larger than the 11.5\% between-school share of variance seen here.

\paragraph{Recommendations for evaluating multi-task models.}
Three practices would have changed the conclusions drawn from School. First, report a pooled model and a partial-pooling model alongside independent models; both are cheap, and the pooled model is the natural null hypothesis when tasks share features. Second, report results over a range of training fractions, since the ranking of methods reverses between 20\% and 30\% (\cref{fig:curves}). Third, state the metric exactly, including how variance is normalized, and check optimizer convergence rather than relying on an iteration budget.

\paragraph{Limitations.}
We study a single benchmark with linear models and squared loss, so we make no claim about other datasets, nonlinear models or classification. Our splits are random per-school splits that follow the published protocols but are not the published splits, which may explain small differences in \cref{tab:repro}. We did not re-implement learned task-relationship methods such as MTRL \citep{zhang2010convex} or GO-MTL \citep{kumar2012learning}; their published numbers suggest gains beyond trace-norm MTL, and combining them with a shared mean is a natural next step. Our hyperparameters are chosen by cross-validation inside each training set, whereas some published results use a separate validation set. Finally, the corrected resampled $t$-test is conservative at small training fractions, so some differences that win on all 10 splits are not marked significant.
